\section{A fixed number of crossings on polynomial time scales}
\label{fixed-sec:crossings}

The polynomial coverage estimate in
Proposition~\ref{poly-prop:coverage} has a direct dynamical application.
Together with the absolute-overlap locking estimate of
Proposition~\ref{static-prop:locking}, it rules out mixing on each fixed
polynomial time scale by using a fixed number of nested crossings.
We give the argument in full. Its number of levels and threshold
intervals remain fixed as $n$ grows.

\begin{theorem}[Polynomial-time obstruction from fixed-depth crossings]
\label{fixed-thm:polynomial}
For every fixed $\beta>1$ and every fixed $a>0$,
\begin{align}
 \PP_{J,\,x_0\sim\pi^J}
 \left\{d^J(x_0,n^a)>\frac14\right\}&\longrightarrow1,
 \label{fixed-eq:continuous-conclusion}\\
 \PP_{J,\,x_0\sim\pi^J}
 \left\{d_{\mathrm{disc}}^J(x_0,\lfloor n^a\rfloor)
                                      >\frac14\right\}
 &\longrightarrow1.
 \label{fixed-eq:discrete-conclusion}
\end{align}
Consequently both
$\PP\{T^J(x_0)>n^a\}$ and
$\PP\{T_{\mathrm{disc}}^J(x_0)>n^a\}$ tend to one under the same joint
law. Each assertion is equivalently a statement that the Gibbs mass of
the indicated initial states tends to one in probability over $J$.
\end{theorem}

As in the main theorem, the transition law in each distance starts from
the realized state $x_0$. The proof uses stationarity to sample possible
targets along a path, and then uses total variation from that fixed
initial state to force a sign passage. Independent random thresholds
will make such a passage unlikely.

\subsection{A stationary path and fixed overlap scales}
\label{fixed-subsec:path-scales}

Fix $\beta>1$ and $a>0$, and abbreviate $\pi^J$ to $\pi$. Put
\begin{equation}\label{fixed-eq:time-parameters}
 t_n=n^a,\qquad m=\lceil2n^{a+1}\rceil,\qquad
 e_*=\frac1{100},\qquad \delta_n=n^{-e_*}.
\end{equation}
Conditional on $J$, start $x_0$ with law $\pi$ and generate an infinite
discrete heat-bath chain $(x_k)_{k\ge0}$. At each step choose a uniform
site and use fresh resampling randomness, including possible
self-transitions. Let $N(t)$ be an independent rate-$n$ Poisson process.
The process $x_{N(t)}$ is the continuous-time chain with rate one at each
site. Indeed, if $K_i^J$ denotes resampling at site $i$, then
$K^J=n^{-1}\sum_iK_i^J$ and its generator is
$\sum_i(K_i^J-I)=n(K^J-I)$. Therefore
\[
 \EE[(K^J)^{N(t)}]=\exp\{nt(K^J-I)\}=P_t^J.
\]
Stationarity gives, for every deterministic $k$,
\begin{equation}\label{fixed-eq:stationary-marginal}
 \mathcal L(x_k\mid J)=\pi.
\end{equation}
Moreover, Chernoff's inequality at parameter $\log2$, with
$\EE N(t_n)=n^{a+1}$ and $m\ge2n^{a+1}$, gives
\begin{equation}\label{fixed-eq:poisson-tail}
 \eta_n:=\PP\{N(t_n)>m\}
 \le\exp\{-(2\log2-1)n^{a+1}\}=o(1).
\end{equation}

Use the common threshold $q_*=q_*(\beta)>0$ of
Proposition~\ref{poly-prop:coverage}. With the same energy cutoff $e_1$
as in the main coverage argument, its choice satisfies
\[
 g=\frac{\beta^2}{4}-\frac{\beta e_1}{2}>0,\qquad
 0<q_*<\min\{1,g/(8\beta^2)\}.
\]
Choose an integer $p\ge2$ such that
\begin{equation}\label{fixed-eq:depth-count}
 pe_*>a+2.
\end{equation}
This integer is fixed throughout the limit. First choose $t_p>0$ with
$2t_p<q_*/2$. For $j=p,p-1,\ldots,2$, take a cutoff
$b_j\in(0,q_*/2)$ allowed by Proposition~\ref{static-prop:locking}
for large-overlap threshold $h=t_j/2$, and then choose $t_{j-1}>0$ with
\begin{equation}\label{fixed-eq:separated-scales}
 t_{j-1}<t_j/10,\qquad 4t_{j-1}<b_j.
\end{equation}
Thus $0<t_1<\cdots<t_p$, and all scale gaps are fixed positive
constants. Independently of the disorder and of all other randomness
introduced below, draw
\begin{equation}\label{fixed-eq:thresholds}
 r_j\sim\operatorname{Unif}[t_j,2t_j],\qquad 1\le j\le p,
\end{equation}
independently for different $j$.

\subsection{Polynomial reference banks cover every possible target}
\label{fixed-subsec:banks}

We next sample finite lists of Gibbs configurations. These lists provide
references for the later crossing construction. Their sizes are chosen
so that coverage holds for every earlier reference choice and every
deterministic index of the path segment $x_0,\ldots,x_m$.

Conditional on $J$, all entries in all banks are independent with law
$\pi$, and all banks are independent of the entire chain. Choose their
sizes recursively. Having chosen $D_{3,i}$ for $i<j$, set
\[
 B_j=\sum_{i<j}(D_{3,i}+1),\qquad B_1=0,
\]
choose $D_{1,j}>B_j+a+3$, and use
Corollary~\ref{poly-cor:coverage-cells} with exponent $D_{1,j}$ to obtain
$D_{3,j}>0$. Bank $j$ consists of
\begin{equation}\label{fixed-eq:bank-sizes}
 M_j=\lceil n^{D_{3,j}+1}\rceil
 \quad\hbox{entries}\quad V_{j,1},\ldots,V_{j,M_j}.
\end{equation}
The constants are chosen before any bank entries or thresholds are
sampled. In particular, the number of tuples with one entry from each
of the first $j-1$ banks satisfies
\begin{equation}\label{fixed-eq:tuple-count}
 Q_j:=\prod_{i<j}M_i\le2^{j-1}n^{B_j}.
\end{equation}
For the empty product, $Q_1=1$.

For such a tuple $\boldsymbol v=(v_1,\ldots,v_{j-1})$, define its cell by
\begin{equation}\label{fixed-eq:cell}
 S_j(\boldsymbol v)
 =\{z\in\Sig: |R(v_i,z)|\le3t_i\text{ for every }i<j\}.
\end{equation}
In particular $S_1=\Sig$. The cell restricts the small overlaps with
the earlier references, while coverage will supply a large overlap
with the current target inside that cell.

\begin{lemma}[Simultaneous bank coverage]
\label{fixed-lem:bank-coverage}
With probability $1-o(1)$, for every $1\le j\le p$, every tuple
$\boldsymbol v$ from banks $1,\ldots,j-1$, and every $0\le k\le m$,
\[
 x_k\in S_j(\boldsymbol v)
 \quad\Longrightarrow\quad
 \text{there is }\ell\le M_j\text{ such that }
 V_{j,\ell}\in S_j(\boldsymbol v),\quad
 |R(V_{j,\ell},x_k)|\ge q_*.
\]
This event depends only on $J$, the chain and the bank lists, and is
independent of the thresholds.
\end{lemma}

\begin{proof}
Intersect the $p$ disorder events from
Corollary~\ref{poly-cor:coverage-cells}. Since $p$ is fixed, the
intersection has probability $1-o(1)$. On it, for every $S\subseteq\Sig$
the set
\[
 \mathcal B_j(S)=
 \{x\in S: \pi\{z\in S:|R(z,x)|\ge q_*\}<n^{-D_{3,j}}\}
\]
has $\pi$-mass less than $n^{-D_{1,j}}$.

Fix a disorder in this intersection and condition on the complete first
$j-1$ banks. Each cell in \eqref{fixed-eq:cell} is then fixed. The path
remains independent of these banks conditional on $J$, so
\eqref{fixed-eq:stationary-marginal} applies to every deterministic
$k$. A union bound over all tuples and all indices gives
\begin{equation}\label{fixed-eq:bad-target-bound}
 \PP\!\left\{
 x_k\in\mathcal B_j(S_j(\boldsymbol v))
 \text{ for some }\boldsymbol v,k
 \,\middle|\, J,\text{ banks }1,\ldots,j-1\right\}
 \le Q_j(m+1)n^{-D_{1,j}}.
\end{equation}
This bound uses marginal stationarity, without requiring independence
between different times.

Now condition also on the entire chain. For a target $x_k$ in its cell
but outside its bad set, acceptable references have $\pi$-mass at least
$n^{-D_{3,j}}$. The entries of bank $j$ are still independent Gibbs
draws under this conditioning. The chance that this bank misses that
set is at most
\[
 (1-n^{-D_{3,j}})^{M_j}
 \le\exp(-M_jn^{-D_{3,j}})\le e^{-n}.
\]
Another union bound gives $Q_j(m+1)e^{-n}$. Combining the two bounds
and summing over the fixed number of banks, the total failure
probability is at most
\begin{equation}\label{fixed-eq:coverage-failure}
 o(1)+\sum_{j=1}^pQ_j(m+1)
              (n^{-D_{1,j}}+e^{-n})=o(1).
\end{equation}
Indeed $Q_j(m+1)=O(n^{B_j+a+1})$ and
$D_{1,j}>B_j+a+3$. Neither part of the event uses any threshold.
\end{proof}

\begin{lemma}[Simultaneous bank locking]
\label{fixed-lem:bank-locking}
With probability $1-o(1)$, for every $i<j\le p$, every entry $v$ of
bank $i$, every entry $w$ of bank $j$, and every $0\le k\le m$,
\begin{equation}\label{fixed-eq:all-bank-locking}
 \begin{gathered}
 |R(x_k,w)|\ge t_j/2,\qquad
 |R(v,x_k)|,\ |R(v,w)|\le b_j
 \\
 \Longrightarrow\qquad
 \bigl||R(v,x_k)|-|R(v,w)|\bigr|\le2\delta_n.
 \end{gathered}
\end{equation}
\end{lemma}

\begin{proof}
For fixed bank indices and deterministic $k$, the triple $(v,x_k,w)$
has law $\pi^{\otimes3}$ conditional on $J$. This follows from the
independence of the two distinct banks and their independence from the
stationary chain. Proposition~\ref{static-prop:locking}, with
$h=t_j/2$ and cutoff $b_j$, bounds the averaged failure probability by
$e^{-n^{9/10}}$. There are only finitely many choices of $j$ for the
large-$n$ threshold in that proposition. A union bound therefore gives
the failure estimate
\begin{equation}\label{fixed-eq:locking-failure}
 \sum_{i<j}(m+1)M_iM_j e^{-n^{9/10}}=o(1).
\end{equation}
Every bank exponent and $p$ are fixed, so the prefactor is polynomial.
\end{proof}

Let $\mathcal G_n$ be the intersection of the events in the two lemmas.
Then
\begin{equation}\label{fixed-eq:good-event}
 \PP(\mathcal G_n^c)=o(1).
\end{equation}
It is determined by $J,x_0,\ldots,x_m$ and the bank lists. We have now
obtained coverage and locking simultaneously for every choice the
crossing construction can make; the thresholds remain independent of
this event.

\subsection{A measurable construction of nested crossings}
\label{fixed-subsec:construction}

For a reference $v\in\Sig$, write $D_v(k)=|R(v,x_k)|$. A single
update changes at most one spin, so
\begin{equation}\label{fixed-eq:one-step}
 |D_v(k+1)-D_v(k)|\le2/n
 \qquad\text{for every }v,k.
\end{equation}
We select the first acceptable entry in each bank, in its given order.
To keep every selected reference defined even on exceptional events,
we specify both successful selections and defaults.

At level $1$, choose $v_1$ as the first entry of bank $1$ satisfying
$D_{v_1}(0)\ge q_*$. If there is no such entry, set $v_1=V_{1,1}$,
$\tau_1=0$, and declare the construction inactive. If selection
succeeds, set
\begin{equation}\label{fixed-eq:first-crossing}
 \tau_1=\min\{0\le k\le m:D_{v_1}(k)\le r_1\}
\end{equation}
when the set is nonempty, and declare the construction active. If that
set is empty, set $\tau_1=0$ and declare it inactive.

Suppose $j\ge2$ and the earlier levels are defined. If the construction
is active and $x_{\tau_{j-1}}\in S_j(v_1,\ldots,v_{j-1})$, select
$v_j$ as the first entry of bank $j$ in that cell with
\[
 |R(v_j,x_{\tau_{j-1}})|\ge q_*.
\]
After a successful selection, search backwards from $\tau_{j-1}$:
\begin{equation}\label{fixed-eq:backward-crossing}
 \tau_j=\max\{0\le k\le\tau_{j-1}:D_{v_j}(k)\le r_j\}
\end{equation}
if this set is nonempty, keeping the construction active. If it is
empty, retain the selected $v_j$, set $\tau_j=0$, and declare the
construction inactive. If the construction was already inactive, the
target was outside its cell, or the reference selection failed, set
$v_j=V_{j,1}$ and $\tau_j=0$, and declare it inactive. An inactive
construction stays inactive at all later levels. At every successful
selection, the depth at the target is at least $q_*>2t_j\ge r_j$;
hence any successful backwards crossing has $\tau_j<\tau_{j-1}$.

The information used in these choices is central to the proof. Define
\[
 \mathcal W=\sigma(J,x_0,\ldots,x_m,\text{ all bank lists}),\qquad
 \mathcal F_0=\mathcal W,\qquad
 \mathcal F_j=\mathcal W\vee\sigma(r_1,\ldots,r_j).
\]
Induction through the explicit rules above gives
\begin{equation}\label{fixed-eq:filtration}
 v_j\text{ is }\mathcal F_{j-1}\text{-measurable},\qquad
 \tau_j\text{ and the active flag after level }j
       \text{ are }\mathcal F_j\text{-measurable}.
\end{equation}
In particular, neither the selection of $v_j$ nor any default for it
uses $r_j$ or a later threshold.

Let $C_1$ denote successful selection of the first reference. On $C_1$
define $s_1=\operatorname{sgn}R(v_1,x_0)\in\{-1,1\}$ and let
\begin{equation}\label{fixed-eq:sign-passage}
 \mathcal S_n=C_1\cap
 \{\text{there exists }0\le k\le m: s_1R(v_1,x_k)\le0\}.
\end{equation}
Take this event to be false outside $C_1$. Starting from signed
overlap at least $q_*$, a sign passage must visit absolute overlap at
most $2/n$, because signed overlaps also change by at most $2/n$.
For sufficiently large $n$, it therefore forces the first crossing
in \eqref{fixed-eq:first-crossing}.

\begin{lemma}[Deterministic nesting]
\label{fixed-lem:nested-crossings}
For all sufficiently large $n$, on $\mathcal G_n\cap\mathcal S_n$
every reference selection and every crossing succeeds. For every
$1\le j\le p$,
\begin{equation}\label{fixed-eq:nested-shells}
 |D_{v_j}(\tau_j)-r_j|\le2/n,\qquad
 |D_{v_i}(\tau_j)-r_i|\le6j\delta_n
                  \quad(1\le i\le j).
\end{equation}
\end{lemma}

\begin{proof}
At level $1$, simultaneous coverage ensures successful selection and
$\mathcal S_n$ ensures a crossing. The initial depth exceeds $r_1$.
Minimality of $\tau_1$ and \eqref{fixed-eq:one-step} therefore put its
depth in $[r_1-2/n,r_1]$, proving the assertions at that level.

Assume the assertions through level $j-1$, where $j\ge2$. For every
$i<j$,
\[
 D_{v_i}(\tau_{j-1})\le2t_i+6(j-1)\delta_n<3t_i
\]
for large $n$. Thus the target is in its cell, and simultaneous coverage
supplies an entry $v_j$ in that cell with new depth at least $q_*$.
Both $D_{v_i}(\tau_{j-1})$ and $|R(v_i,v_j)|$ are at most $3t_i$,
which is less than $b_j$ because $3t_i\le3t_{j-1}<b_j$. The new
depth at the target is at least $q_*>t_j/2$. Simultaneous locking gives
\begin{equation}\label{fixed-eq:reference-depth}
 \bigl||R(v_i,v_j)|-r_i\bigr|
 \le\bigl|D_{v_i}(\tau_{j-1})-r_i\bigr|+2\delta_n
 \le(6(j-1)+2)\delta_n.
\end{equation}

We next verify the locking hypotheses throughout the backwards search.
At any visited index $k$ where the new depth is at least $t_j/2$ and
the old depths are at most $b_j$, locking implies
\[
 D_{v_i}(k)\le |R(v_i,v_j)|+2\delta_n
             \le3t_i+2\delta_n.
\]
The gap $b_j-3t_{j-1}$ is a fixed positive number. Hence
\eqref{fixed-eq:one-step} ensures that every old depth remains below
$b_j$ at the next backwards step, for all sufficiently large $n$.
Before a crossing, the new depth exceeds $r_j\ge t_j$; at the first
crossing step it is at least $r_j-2/n\ge t_j/2$. Starting from the
target, these observations prove by induction over the backwards
steps that locking continues to apply up to and including the first
crossing. If no crossing exists, the induction reaches index $0$.

Throughout that portion of the search, locking and
\eqref{fixed-eq:reference-depth} give, for $i<j$,
\begin{equation}\label{fixed-eq:old-depths}
 |D_{v_i}(k)-r_i|\le(6(j-1)+4)\delta_n.
\end{equation}
A crossing must occur: otherwise this inequality at $k=0$ and $i=1$
would imply
\[
 q_*\le D_{v_1}(0)
 \le2t_1+(6(j-1)+4)\delta_n<q_*.
\]
At the backwards crossing, the next index towards the target has depth
strictly greater than $r_j$. Consequently
$D_{v_j}(\tau_j)\in[r_j-2/n,r_j]$ by the one-step bound.
Combining this with \eqref{fixed-eq:old-depths} proves
\eqref{fixed-eq:nested-shells}. Since $p$, $q_*$, and all $t_i,b_j$ are fixed,
one sufficiently large value of $n$ makes every buffer used in the
induction valid simultaneously.
\end{proof}

\subsection{Independent thresholds make the final shells unlikely}
\label{fixed-subsec:threshold-estimate}

The deterministic construction has reduced a sign passage to a single
time at which all selected reference depths lie in narrow intervals
about their thresholds. Define this event, using the references from
the construction including its defaults, by
\begin{equation}\label{fixed-eq:shell-event}
 \mathcal A_n=
 \{\text{there exists }0\le k\le m:
       |D_{v_j}(k)-r_j|\le6p\delta_n
       \text{ for every }1\le j\le p\}.
\end{equation}
Lemma~\ref{fixed-lem:nested-crossings}, applied at $\tau_p$, shows that
\begin{equation}\label{fixed-eq:sign-implies-shells}
 \mathcal G_n\cap\mathcal S_n\subseteq\mathcal A_n.
\end{equation}

Fix a deterministic $k$ and set
$E_{j,k}=\{|D_{v_j}(k)-r_j|\le6p\delta_n\}$.
By \eqref{fixed-eq:filtration}, its center $D_{v_j}(k)$ is
$\mathcal F_{j-1}$-measurable, while $r_j$ is independent of
$\mathcal F_{j-1}$ and uniform on an interval of length $t_j$.
Thus
\begin{equation}\label{fixed-eq:one-threshold-bound}
 \PP(E_{j,k}\mid\mathcal F_{j-1})
 \le\min\{1,12p\delta_n/t_j\}.
\end{equation}
The earlier events $E_{i,k}$, $i<j$, are
$\mathcal F_{j-1}$-measurable. Applying conditional expectation
successively for $j=p,p-1,\ldots,1$ therefore gives
\[
 \PP\!\left(\bigcap_{j=1}^pE_{j,k}\,\middle|\,\mathcal W\right)
 \le\prod_{j=1}^p\min\{1,12p\delta_n/t_j\}.
\]
The centers may depend on earlier thresholds; the displayed filtration
is exactly what allows this product estimate. A union bound over the
deterministic indices $k=0,\ldots,m$ now yields
\begin{equation}\label{fixed-eq:simultaneous-shell-bound}
 \PP(\mathcal A_n)
 \le(m+1)\prod_{j=1}^p\frac{12p\delta_n}{t_j}
 \le C_{\beta,a}n^{a+1-pe_*}=o(1).
\end{equation}
The exponent is less than $-1$ by \eqref{fixed-eq:depth-count}.

\subsection{A fast realized initial state forces a sign passage}
\label{fixed-subsec:tv-conclusion}

\begin{proof}[Proof of Theorem~\ref{fixed-thm:polynomial}]
First consider continuous time and write
\[
 q_n=\PP_{J,\,x_0\sim\pi^J}
       \{d^J(x_0,t_n)\le1/4\}.
\]
The first reference is chosen using only $J,x_0$ and bank $1$.
Conditional on these variables, the future chain retains its transition
law from the realized $x_0$. On $C_1$, define
\[
 A=A(J,x_0,v_1)=\{z\in\Sig:s_1R(v_1,z)\le0\}.
\]
Spin reversal preserves $\pi$ and interchanges positive and negative
overlaps, so $\pi(A)\ge1/2$. For any realized $J,x_0$ with
$d^J(x_0,t_n)\le1/4$, and any bank $1$ for which $C_1$ holds, the
definition of total variation gives
\[
 \PP\{x_{N(t_n)}\in A\mid J,x_0,\text{bank }1\}\ge1/4.
\]
If $N(t_n)\le m$, this endpoint event implies $\mathcal S_n$.
The Poisson process is independent of all the conditioning variables,
so the error in \eqref{fixed-eq:poisson-tail} is uniform under this
conditioning. Consequently
\begin{align}
 \PP(\mathcal S_n)
 &\ge(1/4-\eta_n)
      \PP\{d^J(x_0,t_n)\le1/4,\ C_1\}\notag\\
 &\ge q_n/4-\eta_n-\tfrac14\PP(C_1^c).
 \label{fixed-eq:fast-start-sign-bound}
\end{align}
The first-bank case of Lemma~\ref{fixed-lem:bank-coverage} gives
$\PP(C_1^c)=o(1)$. Adding the independent thresholds does not change
any probability in this comparison.

Combining \eqref{fixed-eq:good-event},
\eqref{fixed-eq:sign-implies-shells},
\eqref{fixed-eq:simultaneous-shell-bound}, and
\eqref{fixed-eq:fast-start-sign-bound} gives the explicit estimate
\begin{equation}\label{fixed-eq:fast-start-final-bound}
 \frac{q_n}{4}
 \le\eta_n+\tfrac14\PP(C_1^c)+\PP(\mathcal G_n^c)
       +C_{\beta,a}n^{a+1-pe_*}=o(1).
\end{equation}
This proves \eqref{fixed-eq:continuous-conclusion}.

For discrete time set $k_n=\lfloor n^a\rfloor\le m$ and
\[
 q_n^{\mathrm{disc}}
 =\PP_{J,\,x_0\sim\pi^J}
       \{d_{\mathrm{disc}}^J(x_0,k_n)\le1/4\}.
\]
Use the same stationary discrete path, banks, scales and thresholds.
Whenever the realized initial state is in the event defining
$q_n^{\mathrm{disc}}$ and $C_1$ holds, total variation now gives
\[
 \PP\{x_{k_n}\in A\mid J,x_0,\text{bank }1\}\ge1/4.
\]
This endpoint is among $x_0,\ldots,x_m$, so it forces $\mathcal S_n$
with no clock-tail error. The same argument therefore gives
\begin{equation}\label{fixed-eq:discrete-final-bound}
 \frac{q_n^{\mathrm{disc}}}{4}
 \le\tfrac14\PP(C_1^c)+\PP(\mathcal G_n^c)
       +C_{\beta,a}n^{a+1-pe_*}=o(1),
\end{equation}
proving \eqref{fixed-eq:discrete-conclusion} directly.

For each realized finite chain, total variation is nonincreasing in
time, and in continuous time it is also continuous. Hence the first
conclusion implies $T^J(x_0)>n^a$. In discrete time it implies
$T_{\mathrm{disc}}^J(x_0)>\lfloor n^a\rfloor$; this integer inequality
gives $T_{\mathrm{disc}}^J(x_0)>n^a$. Finally, the Gibbs mass of each
exceptional set lies in $[0,1]$ and has expectation tending to zero.
Markov's inequality makes that mass tend to zero in probability over
$J$; conversely, convergence in probability of a bounded mass implies
convergence of its expectation. This proves the asserted equivalent
formulations.
\end{proof}
